\documentclass[../../../include/open-logic-section]{subfiles}

\begin{document}

\olfileid[ar]{sfr}{infinite}{card-sb}

\olsection{ملحق: إثبات مبرهنة شرودر–برنشتاين}

بقي قبل مفارقة نظرية المجموعات الساذجة برهان أخير، ساذج في مبانيه دقيق في صنعته: برهان مبرهنة شرودر–برنشتاين
(\olref[sfr][siz][sb]{thm:schroder-bernstein}). ومفادها أنه إذا كان
$\cardle{{\OLId{latin-looped}{A}}}{{\OLId{latin-looped}{B}}}$ و$\cardle{{\OLId{latin-looped}{B}}}{{\OLId{latin-looped}{A}}}$، كان $\cardeq{{\OLId{latin-looped}{A}}}{{\OLId{latin-looped}{B}}}$؛ أي إنه متى وُجدت
!!{injection}s ${\OLId{latin-ordinary}{f}} \colon {\OLId{latin-looped}{A}} \to {\OLId{latin-looped}{B}}$ و${\OLId{latin-ordinary}{g}} \colon {\OLId{latin-looped}{B}} \to {\OLId{latin-looped}{A}}$، وُجدت
!!a{bijection} ${\OLId{latin-ordinary}{h}} \colon {\OLId{latin-looped}{A}} \to {\OLId{latin-looped}{B}}$.

جرى البحث في هذا الفصل على مفهوم ديديكند في \emph{الإغلاقات}. وبهذا المفهوم أقام ديديكند برهانًا حسنًا لمبرهنة شرودر–برنشتاين، نعرضه الآن. وطريقتنا قريبة من معالجة
\citet[صص.~157--8]{Potter2004}؛ فمن شاء أن يرجع إليها وجد عرضًا يختلف قليلًا ويتفق في الأصل. ويكفي قليل من البحث في غوغل ليتبين أن لهذه المبرهنة براهين
\emph{كثيرة} مختلفة طريفة، كما لعدم نسبية
$\sqrt{{\OLMathNumeral{2}}}$ براهين من هذا القبيل.

وباستخدام ترميز شبيه بما في
\olref[sfr][infinite][dedekind]{Closure}، لتكن ${\OLId{latin-looped}{E}}$ مجموعة، ولتكن
${\OLId{latin-looped}{B}}\subseteq {\OLId{latin-looped}{E}}$ و${\OLId{latin-ordinary}{f}}\colon {\OLId{latin-looped}{E}}\to {\OLId{latin-looped}{E}}$، ونضع
\[
\Closureofunder{{\OLId{latin-ordinary}{f}}}{{\OLId{latin-looped}{B}}} = \bigcap \Setabs{{\OLId{latin-looped}{X}}}{{\OLId{latin-looped}{X}}\subseteq {\OLId{latin-looped}{E}}\text{ و}{\OLId{latin-looped}{B}} \subseteq {\OLId{latin-looped}{X}}
\text{ و${\OLId{latin-looped}{X}}$ منغلقة تحت ${\OLId{latin-ordinary}{f}}$}}
\]
وبحسب هذا التعريف، تكون
$\Closureofunder{{\OLId{latin-ordinary}{f}}}{{\OLId{latin-looped}{B}}}$ أصغر مجموعة منغلقة تحت ${\OLId{latin-ordinary}{f}}$ وتحتوي~${\OLId{latin-looped}{B}}$، بمعنى
ما يأتي:

\begin{lem}\ollabel{Closureprops}
	لكل مجموعة ${\OLId{latin-looped}{E}}$، ولكل دالة ${\OLId{latin-ordinary}{f}}\colon {\OLId{latin-looped}{E}}\to {\OLId{latin-looped}{E}}$، ولكل ${\OLId{latin-looped}{B}}\subseteq {\OLId{latin-looped}{E}}$:
	\begin{enumerate}
		\item\ollabel{Closurehaselem} ${\OLId{latin-looped}{B}} \subseteq \Closureofunder{{\OLId{latin-ordinary}{f}}}{{\OLId{latin-looped}{B}}}$؛ و
		\item\ollabel{Closureclosed} $\Closureofunder{{\OLId{latin-ordinary}{f}}}{{\OLId{latin-looped}{B}}}$ منغلقة تحت ${\OLId{latin-ordinary}{f}}$؛ و
		\item\ollabel{Closuresmallest} إذا كانت ${\OLId{latin-looped}{X}}$ منغلقة تحت ${\OLId{latin-ordinary}{f}}$ وكان ${\OLId{latin-looped}{B}}
		\subseteq {\OLId{latin-looped}{X}}$، فإن $\Closureofunder{{\OLId{latin-ordinary}{f}}}{{\OLId{latin-looped}{B}}} \subseteq {\OLId{latin-looped}{X}}$.
	\end{enumerate}
\end{lem}

\textbf{تنبيه تصحيحي على الأصل (C054-01).} صُرّح بالمجموعة المحيطة ${\OLId{latin-looped}{E}}$ وبكون
${\OLId{latin-ordinary}{f}}$ دالة من ${\OLId{latin-looped}{E}}$ إلى نفسها؛ فبدون ذلك لا يلزم أن يكون تقاطع الجزئيات
منغلقًا تحت ${\OLId{latin-ordinary}{f}}$، ولا يكون المتغير ${\OLId{latin-looped}{X}}$ محصورًا في مجموعة معلومة.

\begin{proof}
يُسلك فيه بعينه مسلك \olref[sfr][infinite][dedekind]{closureproperties}.
\end{proof}

ولا يبقى قبل إثبات شرودر–برنشتاين إلا القضية الآتية:

\begin{prop}\ollabel{sbhelper}
إذا كان ${\OLId{latin-looped}{A}} \subseteq {\OLId{latin-looped}{B}} \subseteq {\OLId{latin-looped}{C}}$ وكان $\cardeq{{\OLId{latin-looped}{A}}}{{\OLId{latin-looped}{C}}}$، فإن $\cardeq{{\OLId{latin-looped}{A}}}{{\OLId{latin-looped}{B}}}$ و$\cardeq{{\OLId{latin-looped}{B}}}{{\OLId{latin-looped}{C}}}$.
\end{prop}

\begin{proof}
إذا أعطيت !!a{bijection} ${\OLId{latin-ordinary}{f}} \colon {\OLId{latin-looped}{C}} \to {\OLId{latin-looped}{A}}$، فبما أن ${\OLId{latin-looped}{A}}\subseteq {\OLId{latin-looped}{C}}$
ننظر إلى ${\OLId{latin-ordinary}{f}}$ أيضًا بوصفها دالة من ${\OLId{latin-looped}{C}}$ إلى نفسها، ولتكن ${\OLId{latin-looped}{F}} =
\Closureofunder{{\OLId{latin-ordinary}{f}}}{{\OLId{latin-looped}{C}} \setminus {\OLId{latin-looped}{B}}}$، وعرّف دالة ${\OLId{latin-ordinary}{g}}$ مجالها ${\OLId{latin-looped}{C}}$ كما
يأتي:
\[
	{\OLId{latin-ordinary}{g}}({\OLId{latin-ordinary}{x}}) =
	\begin{cases}
			{\OLId{latin-ordinary}{f}}({\OLId{latin-ordinary}{x}}) &\text{إذا كان ${\OLId{latin-ordinary}{x}} \in {\OLId{latin-looped}{F}}$}\\
			{\OLId{latin-ordinary}{x}} & \text{فيما عدا ذلك}
		\end{cases}
\]
سنبيّن أن ${\OLId{latin-ordinary}{g}}$ هي !!a{bijection} ${\OLId{latin-looped}{C}} \to {\OLId{latin-looped}{B}}$ (أي إن $\cardeq{{\OLId{latin-looped}{B}}}{{\OLId{latin-looped}{C}}}$)، ومن ذلك يلزم أن
$\comp{{\OLId{latin-ordinary}{f}}^{-{\OLMathNumeral{1}}}}{{\OLId{latin-ordinary}{g}}} \colon {\OLId{latin-looped}{A}} \to {\OLId{latin-looped}{B}}$ هي !!a{bijection} (أي إن $\cardeq{{\OLId{latin-looped}{A}}}{{\OLId{latin-looped}{B}}}$)، فيكتمل البرهان.

نبدأ بتمييز الصورتين: إذا كان ${\OLId{latin-ordinary}{x}} \in {\OLId{latin-looped}{F}}$ وكان ${\OLId{latin-ordinary}{y}}\notin {\OLId{latin-looped}{F}}$، فلا بد أن ${\OLId{latin-ordinary}{g}}({\OLId{latin-ordinary}{x}}) \neq
{\OLId{latin-ordinary}{g}}({\OLId{latin-ordinary}{y}})$. وإلا لزم ${\OLId{latin-ordinary}{y}} = {\OLId{latin-ordinary}{g}}({\OLId{latin-ordinary}{y}}) = {\OLId{latin-ordinary}{g}}({\OLId{latin-ordinary}{x}}) =
{\OLId{latin-ordinary}{f}}({\OLId{latin-ordinary}{x}})$. لكن ${\OLId{latin-ordinary}{x}} \in {\OLId{latin-looped}{F}}$، و${\OLId{latin-looped}{F}}$ منغلقة تحت ${\OLId{latin-ordinary}{f}}$ بمقتضى
\olref{Closureprops}؛ فيكون ${\OLId{latin-ordinary}{y}} = {\OLId{latin-ordinary}{f}}({\OLId{latin-ordinary}{x}}) \in {\OLId{latin-looped}{F}}$، خلاف الفرض.

وعليه، إذا فرضنا ${\OLId{latin-ordinary}{g}}({\OLId{latin-ordinary}{x}}) = {\OLId{latin-ordinary}{g}}({\OLId{latin-ordinary}{y}})$، لزم أن ${\OLId{latin-ordinary}{x}} \in {\OLId{latin-looped}{F}}$ إذا
وفقط إذا كان ${\OLId{latin-ordinary}{y}} \in {\OLId{latin-looped}{F}}$. فإن كان ${\OLId{latin-ordinary}{x}}, {\OLId{latin-ordinary}{y}} \in {\OLId{latin-looped}{F}}$، كان ${\OLId{latin-ordinary}{f}}({\OLId{latin-ordinary}{x}}) = {\OLId{latin-ordinary}{g}} ({\OLId{latin-ordinary}{x}}) =
{\OLId{latin-ordinary}{g}}({\OLId{latin-ordinary}{y}}) = {\OLId{latin-ordinary}{f}}({\OLId{latin-ordinary}{y}})$، فيلزم ${\OLId{latin-ordinary}{x}} = {\OLId{latin-ordinary}{y}}$ لأن ${\OLId{latin-ordinary}{f}}$ هي !!a{bijection}. وإن كان
${\OLId{latin-ordinary}{x}}, {\OLId{latin-ordinary}{y}} \notin {\OLId{latin-looped}{F}}$، حصل ${\OLId{latin-ordinary}{x}} = {\OLId{latin-ordinary}{g}}({\OLId{latin-ordinary}{x}}) = {\OLId{latin-ordinary}{g}}({\OLId{latin-ordinary}{y}}) = {\OLId{latin-ordinary}{y}}$. ففي الحالين يتساوى المدخلان؛ وبذلك تكون ${\OLId{latin-ordinary}{g}}$
!!a{injection}.

ولإتمام البرهان نثبت $\ran{{\OLId{latin-ordinary}{g}}} = {\OLId{latin-looped}{B}}$. خذ ${\OLId{latin-ordinary}{x}} \in {\OLId{latin-looped}{B}} \subseteq {\OLId{latin-looped}{C}}$. إن كان
${\OLId{latin-ordinary}{x}} \notin {\OLId{latin-looped}{F}}$، فلدينا ${\OLId{latin-ordinary}{g}}({\OLId{latin-ordinary}{x}}) = {\OLId{latin-ordinary}{x}}$. وإن كان ${\OLId{latin-ordinary}{x}} \in {\OLId{latin-looped}{F}}$، فلا بد أن ${\OLId{latin-ordinary}{x}} = {\OLId{latin-ordinary}{f}}({\OLId{latin-ordinary}{y}})$
لبعض ${\OLId{latin-ordinary}{y}} \in {\OLId{latin-looped}{F}}$. إذ لو انتفت هذه الصورة، لكانت ${\OLId{latin-looped}{F}} \setminus \{{\OLId{latin-ordinary}{x}}\}$
منغلقة تحت ${\OLId{latin-ordinary}{f}}$ ومحتوية على ${\OLId{latin-looped}{C}}\setminus {\OLId{latin-looped}{B}}$، وهو محال بمقتضى
\olref{Closureprops}. وحينئذ ${\OLId{latin-ordinary}{g}}({\OLId{latin-ordinary}{y}}) = {\OLId{latin-ordinary}{f}}({\OLId{latin-ordinary}{y}}) = {\OLId{latin-ordinary}{x}}$.
\end{proof}

فلنرجع إلى المبرهنة الأصلية. نذكّر بأن صورة الدالة ${\OLId{latin-ordinary}{h}}$ على المجموعة ${\OLId{latin-looped}{D}}$ معرّفة هكذا:
$\funimage{{\OLId{latin-ordinary}{h}}}{{\OLId{latin-looped}{D}}} = \Setabs{{\OLId{latin-ordinary}{h}}({\OLId{latin-ordinary}{x}})}{{\OLId{latin-ordinary}{x}}
\in {\OLId{latin-looped}{D}}}$.

\begin{proof}[برهان مبرهنة شرودر–برنشتاين] خذ !!{injection}s
${\OLId{latin-ordinary}{f}} \colon {\OLId{latin-looped}{A}} \to {\OLId{latin-looped}{B}}$ و${\OLId{latin-ordinary}{g}} \colon {\OLId{latin-looped}{B}} \to {\OLId{latin-looped}{A}}$. من $\funimage{{\OLId{latin-ordinary}{f}}}{{\OLId{latin-looped}{A}}}
\subseteq {\OLId{latin-looped}{B}}$ يلزم $\funimage{{\OLId{latin-ordinary}{g}}}{\funimage{{\OLId{latin-ordinary}{f}}}{{\OLId{latin-looped}{A}}}} \subseteq
{\OLId{latin-ordinary}{g}}[{\OLId{latin-looped}{B}}] \subseteq {\OLId{latin-looped}{A}}$. والدالة $\comp{{\OLId{latin-ordinary}{f}}}{{\OLId{latin-ordinary}{g}}} \colon {\OLId{latin-looped}{A}} \to
\funimage{{\OLId{latin-ordinary}{g}}}{\funimage{{\OLId{latin-ordinary}{f}}}{{\OLId{latin-looped}{A}}}}$ هي !!{injection}، إذ كل من ${\OLId{latin-ordinary}{g}}$ و
${\OLId{latin-ordinary}{f}}$ كذلك؛ ثم إن $\comp{{\OLId{latin-ordinary}{f}}}{{\OLId{latin-ordinary}{g}}}$ هي !!a{bijection} أيضًا، لاختيار صورتها مجالًا مقابلًا. فينتج
$\cardeq{\funimage{{\OLId{latin-ordinary}{g}}}{\funimage{{\OLId{latin-ordinary}{f}}}{{\OLId{latin-looped}{A}}}}}{{\OLId{latin-looped}{A}}}$، ومن
\olref{sbhelper} توجد !!a{bijection} ${\OLId{latin-ordinary}{h}} \colon {\OLId{latin-looped}{A}} \to
\funimage{{\OLId{latin-ordinary}{g}}}{{\OLId{latin-looped}{B}}}$. وأما ${\OLId{latin-ordinary}{g}}^{-{\OLMathNumeral{1}}}$ فهي !!a{bijection}
$\funimage{{\OLId{latin-ordinary}{g}}}{{\OLId{latin-looped}{B}}} \to {\OLId{latin-looped}{B}}$. فتركيبهما $\comp{{\OLId{latin-ordinary}{h}}}{{\OLId{latin-ordinary}{g}}^{-{\OLMathNumeral{1}}}} \colon {\OLId{latin-looped}{A}} \to {\OLId{latin-looped}{B}}$
!!a{bijection}، وهو المطلوب.
\end{proof}

\end{document}
